\documentclass[a4paper,11pt]{article}

% --- Packages fondamentaux ---
\usepackage[utf8]{inputenc}
\usepackage[T1]{fontenc}
\usepackage[french]{babel}
\usepackage{geometry}
\geometry{top=2cm, bottom=2cm, left=2.2cm, right=2.2cm}
\usepackage{amsmath, amssymb}
\usepackage{graphicx}
\usepackage{booktabs}
\usepackage{hyperref}
\usepackage{xcolor}
\usepackage{fancyhdr}
\usepackage{tcolorbox}
\tcbuselibrary{skins}
\usepackage{titlesec}

% --- Couleurs ---
\definecolor{primary}{RGB}{20, 50, 110}
\definecolor{secondary}{RGB}{40, 120, 180}
\definecolor{darkgray}{RGB}{50, 50, 50}
\definecolor{codeblue}{RGB}{20, 70, 180}

% --- Hyperliens ---
\hypersetup{
    colorlinks=true,
    linkcolor=primary,
    filecolor=primary,
    urlcolor=codeblue,
    citecolor=primary,
    pdftitle={Rapport Devoir 2 - Transformation de Perspective - Yahya AHMANE},
    pdfauthor={Yahya AHMANE}
}

% --- En-têtes et pieds de page ---
\pagestyle{fancy}
\fancyhf{}
\fancyhead[L]{\small\color{gray}FST Tanger --- LSI | Vision Artificielle}
\fancyhead[R]{\small\color{gray}Devoir 2 : Transformation de Perspective}
\fancyfoot[C]{\thepage}
\renewcommand{\headrulewidth}{0.4pt}

% --- Formatage des titres ---
\titleformat{\section}{\color{primary}\large\bfseries}{\thesection.}{0.6em}{}[\titlerule]
\titleformat{\subsection}{\color{secondary}\normalsize\bfseries}{\thesubsection.}{0.5em}{}

\begin{document}

% =============================================================
%                       PAGE DE GARDE
% =============================================================
\begin{titlepage}
    \centering
    \vspace*{0.5cm}
    
    {\scshape\Large Université Abdelmalek Essaâdi \par}
    {\scshape\large Faculté des Sciences et Techniques de Tanger \par}
    {\scshape\footnotesize Département de Génie Informatique \par}
    
    \vspace{0.8cm}
    \hrule height 1.5pt
    \vspace{0.8cm}
    
    {\large\bfseries Cycle d'Ingénieur en Logiciels et Systèmes Intelligents (LSI) \par}
    \vspace{0.3cm}
    {\large\color{secondary}\textbf{Module : Vision Artificielle \& Traitement d'Images} \par}
    
    \vspace{2.0cm}
    
    \begin{tcolorbox}[colback=primary!5!white,colframe=primary,arc=3mm,boxrule=1.2pt]
        \centering
        \vspace{0.3cm}
        {\LARGE\bfseries\color{primary} DEVOIR 2 : ALIGNEMENT ET TRANSFORMATION DE PERSPECTIVE D'UN DOCUMENT \par}
        \vspace{0.3cm}
        {\large Orthorectification Automatique et Interactive avec OpenCV (C++) \par}
        \vspace{0.2cm}
    \end{tcolorbox}
    
    \vspace{2.5cm}
    
    \begin{minipage}[t]{0.48\textwidth}
        \textbf{\large Réalisé par :}\\[0.2cm]
        {\Large\textbf{Yahya AHMANE}}\\[0.1cm]
        Élève Ingénieur en LSI\\
        \href{mailto:ahmaneyahya10@gmail.com}{\texttt{ahmaneyahya10@gmail.com}}
    \end{minipage}
    \hfill
    \begin{minipage}[t]{0.48\textwidth}
        \begin{flushright}
            \textbf{\large Encadré par :}\\[0.2cm]
            {\Large\textbf{Pr. M'hamed AIT KBIR}}\\[0.1cm]
            Professeur d'Enseignement Supérieur\\
            FST de Tanger
        \end{flushright}
    \end{minipage}
    
    \vfill
    
    \begin{tcolorbox}[colback=gray!10!white,colframe=gray!40!white,arc=2mm,boxrule=0.5pt,center]
        \centering\footnotesize
        \textbf{Dépôt GitHub du Code Source :} \href{https://github.com/ahyahya1616/transformation_perspective}{\texttt{https://github.com/ahyahya1616/transformation_perspective}}
    \end{tcolorbox}
    
    \vspace{0.5cm}
    {\small Année Universitaire 2024 -- 2025 \par}
\end{titlepage}

% =============================================================
% 1. CONTEXTE ET OBJECTIF DU TRAVAIL
% =============================================================
\section{Contexte et Objectif Général}

Lors de la prise de vue photographique d'un document posé sur un support, l'inclinaison de l'appareil par rapport au plan du document engendre une déformation perspective. L'objectif de ce devoir est d'écrire un programme en \textbf{C++} utilisant la bibliothèque \textbf{OpenCV} permettant de sélectionner le quadrilatère du document et de lui appliquer une transformation de perspective afin d'obtenir une vue de dessus redressée (orthorectifiée).

Le programme intègre les 4 étapes exigées :
\begin{enumerate}
    \item Une détection automatique personnalisée des 4 coins du document ;
    \item Un affichage graphique du résultat de cette détection ;
    \item Une phase d'ajustement interactif des points détectés ;
    \item L'application mathématique de la transformation de perspective.
\end{enumerate}

Le système prend en charge deux types d'images : les images à \textbf{fond simple} (uniforme et contrasté) et les images à \textbf{fond texturé} (complexe ou bruité).

% =============================================================
% 2. DETECTION AUTOMATIQUE DES 4 COINS ET JUSTIFICATIONS
% =============================================================
\section{Détection Automatique des 4 Coins et Choix des Méthodes}

\subsection{Cas des Images à Fond Simple (Uniforme et Contrasté)}

\paragraph{Justification du choix :}
Dans le cas d'un fond uniforme et contrasté (par exemple un reçu blanc sur une table foncée), la distribution des intensités lumineuses est nettement bimodale. Le document forme une composante connexe homogène avec une frontière nette. La méthode basée sur la binarisation globale d'Otsu et l'extraction de contours (\texttt{findContours} suivi de l'approximation polygonale \texttt{approxPolyDP}) est le choix optimal : elle est directe, déterministe, ultra-rapide et isole immédiatement le quadrilatère de surface maximale.

\paragraph{Déroulement de l'algorithme :}
\begin{enumerate}
    \item \textbf{Conversion et lissage :} Conversion en niveaux de gris et application d'un flou gaussien ($5 \times 5$) pour éliminer les bruits de capteur.
    \item \textbf{Seuillage d'Otsu :} Calcul automatique du seuil optimal séparant le fond de la page blanche pour obtenir un masque binaire.
    \item \textbf{Fermeture morphologique :} Application d'une opération de fermeture (dilatation puis érosion avec un élément structurant $7 \times 7$) pour combler les zones de texte à l'intérieur de la feuille.
    \item \textbf{Extraction et approximation polygonale :} Extraction du contour externe d'aire maximale avec \texttt{cv::findContours}, puis simplification par l'algorithme de Ramer-Douglas-Peucker (\texttt{cv::approxPolyDP} avec une tolérance $\epsilon = 0.02 \times \text{périmètre}$) pour obtenir un polygone fermé à 4 sommets.
\end{enumerate}

\subsection{Cas des Images à Fond Texturé (Complexe ou Bruité)}

\paragraph{Justification du choix :}
Sur un fond texturé (veines de bois marquées, tissu plissé, dégradés d'ombres), la binarisation par seuillage échoue car les intensités des pixels du fond chevauchent celles du papier, créant des contours morcelés. Pour surmonter cette difficulté, nous exploitons un invariant géométrique fort : \textbf{la rectilinéarité des 4 arêtes extérieures du document}. En combinant un filtre bilatéral (qui lisse la texture tout en préservant la netteté des bords), le détecteur de Canny et la Transformée de Hough Probabiliste (\texttt{HoughLinesP}), nous extrayons les lignes maîtresses du document et calculons leurs points d'intersection pour retrouver les 4 coins, même en présence d'un fond complexe ou d'occlusions locales.

\paragraph{Déroulement de l'algorithme :}
\begin{enumerate}
    \item \textbf{Filtrage bilatéral :} Lissage des micro-textures de l'arrière-plan sans flouter les arêtes du document grâce à une pondération spatiale et photométrique.
    \item \textbf{Détection de contours par Canny :} Extraction des gradients fins avec un seuillage par hystérésis adapté ($T_{low}=20, T_{high}=80$).
    \item \textbf{Transformée de Hough Probabiliste (\texttt{HoughLinesP}) :} Détection des segments de droites d'une longueur minimale pour ne retenir que les segments structurants.
    \item \textbf{Regroupement et fusion des droites :} Séparation des segments détectés selon leur inclinaison angulaire en deux catégories : bords quasi-horizontaux ($\alpha \le 25^\circ$) et quasi-verticaux ($\alpha \ge 65^\circ$), puis fusion des segments colinéaires.
    \item \textbf{Calcul géométrique des 4 coins :} Calcul précis des intersections entre les droites maîtresses (haut/gauche, haut/droite, bas/droite, bas/gauche) à l'aide des déterminants vectoriels.
\end{enumerate}

% =============================================================
% 3. AFFICHAGE DU RESULTAT DE LA DETECTION
% =============================================================
\section{Affichage du Résultat de la Détection}

Après détection automatique, les 4 coins sont ordonnés spatialement dans le sens horaire :
\begin{equation*}
    0 : \text{Haut-Gauche (TL)}, \quad 1 : \text{Haut-Droite (TR)}, \quad 2 : \text{Bas-Droite (BR)}, \quad 3 : \text{Bas-Gauche (BL)}
\end{equation*}
L'affichage présente :
\begin{itemize}
    \item Des cercles colorés centrés sur chaque coin avec leur numéro et libellé respectifs ;
    \item Un quadrilatère vert traçant le contour fermé reliant les 4 points ;
    \item Un bandeau supérieur affichant les instructions d'aide pour l'utilisateur.
\end{itemize}

% =============================================================
% 4. PHASE D'AJUSTEMENT INTERACTIF
% =============================================================
\section{Phase d'Ajustement Interactif des Points Détectés}

Pour permettre à l'utilisateur d'affiner manuellement le cadrage ou de corriger une légère imprécision :
\begin{itemize}
    \item \textbf{Interaction à la souris :} Grâce à \texttt{cv::setMouseCallback}, un clic gauche sur un coin le sélectionne, le déplacement de la souris repositionne le point en temps réel, et le relâchement confirme la nouvelle position.
    \item \textbf{Zone de débordement (\textit{Padding}) :} Pour manipuler aisément des coins situés à l'extrême limite du cadre (comme dans \texttt{cell\_pic.jpg}), une marge neutre grise de $120$ pixels est ajoutée autour de l'image (\texttt{copyMakeBorder}), évitant tout blocage au bord de la fenêtre.
    \item \textbf{Contrôles clavier :}
    \begin{itemize}
        \item \texttt{[ENTRÉE]} : Valide les 4 sommets et applique l'orthorectification ;
        \item \texttt{[R]} : Réinitialise les sommets sur la détection automatique initiale ;
        \item \texttt{[Q]} : Quitte le programme.
    \end{itemize}
\end{itemize}

% =============================================================
% 5. TRANSFORMATION DE PERSPECTIVE
% =============================================================
\section{Application de la Transformation de Perspective}

La transformation projective (homographie 2D) est calculée à partir des 4 points sources détectés $P_{src} = \{TL, TR, BR, BL\}$ et des 4 points cibles du rectangle redressé $P_{dst} = \{(0,0), (W-1,0), (W-1,H-1), (0,H-1)\}$.

\begin{enumerate}
    \item \textbf{Dimensions de l'image redressée :} Pour préserver les proportions réelles sans anamorphose, la largeur et la hauteur cibles sont calculées par la distance euclidienne maximale entre coins opposés :
    \begin{equation}
        W = \max\left(\|TR - TL\|, \, \|BR - BL\|\right), \quad
        H = \max\left(\|BL - TL\|, \, \|BR - TR\|\right)
    \end{equation}
    \item \textbf{Calcul de la matrice d'homographie :} Appel de \texttt{cv::getPerspectiveTransform(src, dst)} qui résout le système linéaire de 8 équations associé à la matrice $3 \times 3$.
    \item \textbf{Orthorectification :} Application de \texttt{cv::warpPerspective(src, dst, H, Size(W, H))} avec interpolation bilinéaire, produisant l'image finale redressée qui est affichée et enregistrée automatiquement sur disque.
\end{enumerate}

% =============================================================
% 6. TESTS, VALIDATION ET AJOUT D'UNE IMAGE DE TEST
% =============================================================
\section{Tests, Validation et Ajout d'une Image de Test Spécifique}

\subsection{Images Fournies par le Sujet}
Le programme a été testé sur l'ensemble des images fournies en pièce jointe :
\begin{itemize}
    \item \texttt{receipt.jpg} : Reçu sur fond noir uniforme $\to$ Détection parfaite via la méthode pour fond simple (Otsu).
    \item \texttt{cell\_pic.jpg} : Feuille de texte photographiée occupant presque toute la largeur $\to$ Détection complète sans confusion avec le cadre du capteur.
    \item \texttt{desk.JPG} : Document sur un bureau en bois clair $\to$ Les contrastes de niveau de gris permettent une détection rapide et fidèle.
\end{itemize}

\subsection{Ajout d'une Image de Test pour le Fond Texturé (\texttt{test\_texture\_2.jpg})}
Les images initialement fournies correspondant principalement à des contrastes francs ou des fonds relativement uniformes, \textbf{nous avons expressément ajouté une image supplémentaire de test (\texttt{test\_texture\_2.jpg})} dans le dossier \texttt{images/}. 

Cette image présente un carnet oblique posé sur un tissu en soie beige avec des plis profonds, des reflets lumineux et un trombone métallique sur le bord supérieur. Ce test valide de façon rigoureuse notre algorithme pour \textbf{fond texturé} :
\begin{itemize}
    \item La méthode d'Otsu échoue sur cette image à cause des plis de la soie ;
    \item Notre méthode par \textbf{Filtrage bilatéral + Canny + Hough + Intersections} réussit parfaitement à isoler les 4 bords directeurs et à calculer les 4 coins exacts du carnet.
\end{itemize}

\begin{table}[h!]
\centering
\small
\begin{tabular}{@{}llcc@{}}
\toprule
\textbf{Image de Test} & \textbf{Type de Fond} & \textbf{Méthode Utilisée} & \textbf{Résultat Détection} \\ \midrule
\texttt{receipt.jpg} & Fond simple (sombre) & Otsu / Contours & \textbf{100\% automatique} \\
\texttt{cell\_pic.jpg} & Fond simple (bois) & Canny / Contours & \textbf{100\% automatique} \\
\texttt{desk.JPG} & Fond simple (bureau) & Canny / Contours & \textbf{100\% automatique} \\
\textbf{\texttt{test\_texture\_2.jpg}} (\textit{ajoutée}) & \textbf{Fond texturé (soie plissée)} & \textbf{Hough / Intersections} & \textbf{100\% automatique} \\
\bottomrule
\end{tabular}
\caption{Synthèse des résultats de détection automatique selon la typologie de fond.}
\end{table}

% =============================================================
% 7. CONCLUSION
% =============================================================
\section{Conclusion}

Ce projet remplit l'intégralité des fonctionnalités demandées :
\begin{itemize}
    \item La détection automatique est adaptée et justifiée selon la nature du fond (contours pour fond simple, lignes de Hough pour fond texturé) ;
    \item L'affichage graphique et l'IHM interactive offrent un contrôle total à l'utilisateur ;
    \item La transformation de perspective restitue fidèlement un document redressé prêt à la lecture.
\end{itemize}

\vspace{0.4cm}
\begin{tcolorbox}[colback=primary!5!white,colframe=primary,arc=2mm,boxrule=0.8pt]
    \centering\small
    \textbf{Code source complet et instructions de compilation (CMake / Makefile) :} \\
    \url{https://github.com/ahyahya1616/transformation_perspective}
\end{tcolorbox}

\end{document}
